\dish{Original Plum Torte}
\altdish{Plum Torte, New York Times Original}
\altdish{Tarte aux prunes, New York Times Original}
\serves{6}
%\makes{}
\prep{75 minutes}
\source{nytimes}

A recipe by Marian Burros, related by Andrew Scrivani for the
\textit{New York Times}.

\begin{ingredients}
  \ingr{}{}{}
  \ingr{}{}{}
  \ingr{}{}{}
  \ingr{}{}{}
  \ingr{}{}{}
  \ingr{}{}{}
  \ingr{}{}{}
  \ingr{}{}{}
  \ingr{}{}{}
  \ingr{}{}{}
\end{ingredients}


\begin{recipe}
  \begin{enumerate}

  \item 

  \end{enumerate}
\end{recipe}

%\accord{}


Original Plum Torte

By Marian Burros

Original Plum Torte

Andrew Scrivani for The New York Times

Time
    1 hour 15 minutes
Rating
    5(10533)
Notes
    Read community notes

The Times published Marian Burros’s recipe for Plum Torte every September from 1983 until 1989, when the
editors determined that enough was enough. The recipe was to be printed for the last time that year. “To
counter anticipated protests,” Ms. Burros wrote a few years later, “the recipe was printed in larger
type than usual with a broken-line border around it to encourage clipping.” It didn’t help. The paper
was flooded with angry letters. “The appearance of the recipe, like the torte itself, is bittersweet,”
wrote a reader in Tarrytown, N.Y. “Summer is leaving, fall is coming. That's what your annual recipe is
all about. Don't be grumpy about it.” We are not! And we pledge that every year, as summer gives way to
fall, we will make sure that the recipe is easily available to one and all. The original 1983 recipe
called for 1 cup sugar; the 1989 version reduced that to ¾ cup. We give both options below. Here are
five ways to adapt the torte.

Ingredients

Yield:8 servings

  • ¾ to 1cup sugar
  • ½cup unsalted butter, softened
  • 1cup unbleached flour, sifted
  • 1teaspoon baking powder
  • Pinch of salt (optional)
  • 2eggs
  • 24halves pitted purple plums
  • Sugar, lemon juice and cinnamon, for topping

Add to Your Grocery List
Ingredient Substitution Guide
Nutritional Information

Nutritional analysis per serving (8 servings)

364 calories; 13 grams fat; 8 grams saturated fat; 0 grams trans fat; 4 grams monounsaturated fat; 1
gram polyunsaturated fat; 60 grams carbohydrates; 4 grams dietary fiber; 45 grams sugars; 5 grams
protein; 81 milligrams sodium

Note: The information shown is Edamam’s estimate based on available ingredients and preparation. It
should not be considered a substitute for a professional nutritionist’s advice.

Powered by
[image]

Preparation

 1. Step 1

    Heat oven to 350 degrees.

 2. Step 2

    Cream the sugar and butter in a bowl. Add the flour, baking powder, salt and eggs and beat well.

 3. Step 3

    Spoon the batter into a springform pan of 8, 9 or 10 inches. Place the plum halves skin side up on
    top of the batter. Sprinkle lightly with sugar and lemon juice, depending on the sweetness of the
    fruit. Sprinkle with about 1 teaspoon of cinnamon, depending on how much you like cinnamon.

 4. Step 4

    Bake 1 hour, approximately. Remove and cool; refrigerate or freeze if desired. Or cool to lukewarm
    and serve plain or with whipped cream. (To serve a torte that was frozen, defrost and reheat it
    briefly at 300 degrees.)

Tip

  • To freeze, double-wrap the torte in foil, place in a plastic bag and seal.